\chapter{Conclusion}
\label{chap:conclusion}

In 2018, the four of us built, in a few weeks, a segmentation that the researchers approved, and a
viewer for it. We had identified the real difficulty, telling vessels from bone, and answered it with
morphology tuned by eye.

Eight years later, pyBones keeps the goal and changes the method. The main lessons I draw from the
rewrite are:
\begin{itemize}
  \item \textbf{Measure before improving.} Without a ground truth, a synthetic phantom calibrated on the
        real scan made every change measurable, and showed me that an intuitive improvement (the first 3D
        refinement) made things worse until the vessel edges were right.
  \item \textbf{Physics over parameters.} The largest gains came from modelling how the image forms
        (partial volume, blur, density) rather than from tuning radii: half-maximum edges for vessels
        and bursts, physical units for every size.
  \item \textbf{3D where 2D cannot decide.} Connectivity and shape in 3D resolve most of the blobs that a
        single slice leaves ambiguous; what remains is a property of the scan (resolution) and of the
        anatomy (rods are tubes too).
\end{itemize}

\section*{Future work}
\addcontentsline{toc}{section}{Future work}

\begin{itemize}
  \item \textbf{Quantification}, the original goal: vessel centrelines, network graph, node and branch
        counts, calibre (with the resolution caveat of Section~\ref{sec:thin}) and branch lengths.
  \item \textbf{Gap bridging} along the vessel direction, to reconnect vessels interrupted by clots.
  \item \textbf{Expert validation} of a few real crops with the annotation page.
  \item \textbf{A web page} on top of the web server: slices with their labels, and an in-browser 3D
        view of lighter meshes.
\end{itemize}
